\subsection{Flat modules}

DEFINITION 1. --- The A-module E is said to be flat if, for every exact sequence of right A-modules and homomorphisms

\[
\mathbf {M} ^ {\prime} \xrightarrow {u} \mathbf {M} \xrightarrow {v} \mathbf {M} ^ {\prime \prime},\tag{11}
\]

the sequence of Z-linear maps

\[
\mathbf {M} ^ {\prime} \otimes_ {\mathbf {A}} \mathrm{E} \xrightarrow {u \otimes 1} \mathbf {M} \otimes_ {\mathbf {A}} \mathrm{E} \xrightarrow {v \otimes 1} \mathbf {M} ^ {\prime \prime} \otimes_ {\mathbf {A}} \mathrm{E}\tag{12}
\]

is exact.

PROPOSITION 3. --- For the A-module E to be flat, it is necessary and sufficient that, for every injective A-homomorphism \(u: \mathbf{M}' \to \mathbf{M}\) of right A-modules, the map \(u \otimes 1: \mathbf{M}' \otimes_{\mathbf{A}} \mathbf{E} \to \mathbf{M} \otimes_{\mathbf{A}} \mathbf{E}\) be injective.

If E is flat and \(u: \mathbf{M}' \to \mathbf{M}\) injective, the sequence \(0 \longrightarrow \mathbf{M}' \stackrel{u}{\longrightarrow} \mathbf{M}\) is exact, hence also

the sequence 0 \(\longrightarrow\) M' \(\otimes_{\mathbf{A}}\) E \(\xrightarrow{u\otimes1}\) M \(\otimes_{\mathbf{A}}\) E, and \(u\otimes1\) is injective. Conversely, consider the exact sequence (11); set M'\,' \(_{1}\) = v(M), and let i : M'\,' \(_{1}\) → M'\,' be the canonical injection

and \(p:\mathbf{M}\to\mathbf{M}_{1}^{\prime\prime}\) the map \(m\mapsto v(m)\) . The sequence \(\mathbf{M}^{\prime}\xrightarrow{u}\mathbf{M}\xrightarrow{p}\mathbf{M}_{1}^{\prime\prime}\longrightarrow0\)

is exact; by II, p.~58, prop. 5, the sequence \(\mathbf{M}' \otimes_{\mathbf{A}} \mathbf{E} \xrightarrow{u \otimes 1} \mathbf{M} \otimes_{\mathbf{A}} \mathbf{E} \xrightarrow{p \otimes 1} \mathbf{M}_{1}'' \otimes_{\mathbf{A}} \mathbf{E}\) is therefore exact. Moreover, we have \(v = i \circ p\) , hence \((v \otimes 1) = (i \otimes 1) \circ (p \otimes 1)\) ; if E satisfies the condition of the statement, then \(i \otimes 1\) is injective, therefore

\[
\operatorname{Ker} (v \otimes 1) = \operatorname{Ker} (p \otimes 1) = \operatorname{Im} (u \otimes 1)
\]

and the sequence (12) is exact.

PROPOSITION 4. ---(i) Let \((\mathbf{E}_i)_{i\in \mathbf{I}}\) be a family of A-modules, \(\mathrm{E} = \bigoplus_{i\in \mathbf{I}}\mathrm{E}_i\) their direct sum

For the A-module E to be flat, it is necessary and sufficient that each of the \(E_{i}\) be so.

\begin{enumerate}
\def\labelenumi{(\roman{enumi})}
\setcounter{enumi}{1}
\tightlist
\item
  Let I be a right-filtered preordered set, \((\mathrm{E}_{\alpha}, f_{\beta\alpha})\) an inductive system of A-modules relative to I, \(E = \varinjlim E_{\alpha}\) its inductive limit. If each of the A-modules \(E_{\alpha}\) is flat, then E is flat.
\end{enumerate}

Let \(M^{\prime} \rightarrow M \rightarrow M^{\prime\prime}\) be an exact sequence of right A-modules.

\begin{enumerate}
\def\labelenumi{(\roman{enumi})}
\item
  For the sequence \(\bigoplus_{i\in I}(\mathbf{M}^{\prime}\otimes_{\mathbf{A}}\mathbf{E}_{i})\to \bigoplus_{i\in I}(\mathbf{M}\otimes_{\mathbf{A}}\mathbf{E}_{i})\to \bigoplus_{i\in I}(\mathbf{M}^{\prime \prime}\otimes_{\mathbf{A}}\mathbf{E}_{i})\) to be exact, it is necessary and sufficient that each of the sequences \(\mathbf{M}^{\prime}\otimes_{\mathbf{A}}\mathbf{E}_{i}\to \mathbf{M}\otimes_{\mathbf{A}}\mathbf{E}_{i}\to \mathbf{M}^{\prime \prime}\otimes_{\mathbf{A}}\mathbf{E}_{i}\) be exact (II, p.~13, prop. 7), which proves (i) since \(\oplus (\mathbf{M}\otimes_{\mathbf{A}}\mathbf{E}_{i})\) is canonically identified with \(\mathbf{M}\otimes_{\mathbf{A}}\mathbf{E}\) (II, p.~61, prop. 7).
\item
  By hypothesis, each of the sequences \(\mathbf{M}' \otimes_{\mathbf{A}} \mathbf{E}_i \to \mathbf{M} \otimes_{\mathbf{A}} \mathbf{E}_i \to \dot{\mathbf{M}}'' \otimes_{\mathbf{A}} \mathbf{E}_i\) is exact, hence so is the sequence \(\mathbf{M}' \otimes_{\mathbf{A}} \mathbf{E} \to \mathbf{M} \otimes_{\mathbf{A}} \mathbf{E} \to \mathbf{M}'' \otimes_{\mathbf{A}} \mathbf{E}\) , since passage to the inductive limit commutes with tensor product (II, p.~93, prop. 7) and preserves exactness (II, p.~91, prop. 3).
\end{enumerate}

Examples. --- 1) It is clear that \(A_{s}\) is a flat A-module; it then follows from prop. 4 (i) that every free A-module, and more generally every projective A-module, is flat (see also II, p.~63, cor. 6).

* Conversely, every flat A-module of finite presentation is projective (no. 5). *

\begin{enumerate}
\def\labelenumi{\arabic{enumi})}
\setcounter{enumi}{1}
\item
  By prop. 4 (ii), every A-module which is the inductive limit of a filtered inductive system of free A-modules is flat. We shall prove a converse in no. 6.
\item
  If A is semisimple, every A-module is projective (VIII, § 5, no. 1, prop. 1), hence flat.
\item
  * If A is an artinian local ring (not necessarily commutative), an A-module is flat if and only if it is free (AC II, § 3, no. 2, cor. 2 of prop. 5). *
\item
  If A is an integral domain, the field of fractions K of A is a flat A-module (II, p.~118, prop. 27).
\item
  * In AC II and III, we shall study two important examples of flat A-modules when A is commutative: the rings of fractions \(S^{-1}\) A, and when A is noetherian, the separated completions of A for the J-adic topologies. *
\item
  Let \(a \in \mathbf{A}\) such that the map \(a_{\mathbf{A}}: x \mapsto ax\) of \(\mathbf{A}\) to \(\mathbf{A}\) is injective (“ \(a\) is not a left zero-divisor”). If \(\mathbf{E}\) is a flat \(\mathbf{A}\) -module, then the homothety \(a_{\mathbf{E}}\) is injective, since it is identified with \(a_{\mathbf{A}} \otimes 1: \mathbf{A}_{d} \otimes_{\mathbf{A}} \mathbf{E} \to \mathbf{A}_{d} \otimes_{\mathbf{A}} \mathbf{E}\) . In particular, if \(\mathbf{A}\) is an integral domain, every flat \(\mathbf{A}\) -module is torsion-free. Conversely, if \(\mathbf{A}\) is principal, every torsion-free \(\mathbf{A}\) -module is flat: indeed, if the \(\mathbf{A}\) -module \(\mathbf{E}\) is torsion-free, every finitely generated submodule of \(\mathbf{E}\) is free (VII, § 4, no. 4, cor. 2 to th. 4), and \(\mathbf{E}\) is a filtering increasing union of flat submodules, hence is flat (prop. 4 (ii)).
\item
  Let B be a ring and \(\rho: \mathbf{A} \to \mathbf{B}\) a homomorphism. If E is a flat A-module, the B-module \(\mathbf{E}_{(\mathbf{B})} = \mathbf{B} \otimes_{\mathbf{A}} \mathbf{E}\) is flat. Indeed, let \(u: \mathbf{N}' \to \mathbf{N}\) be an injective homomorphism of right B-modules; then \(u \otimes_{\mathbf{B}} 1_{\mathbf{E}_{(\mathbf{B})}}\) is canonically identified with the homomorphism \(u \otimes_{\mathbf{A}} 1_{\mathbf{E}}: \mathbf{N}' \otimes_{\mathbf{A}} \mathbf{E} \to \mathbf{N} \otimes_{\mathbf{A}} \mathbf{E}\) , which is injective if E is flat.
\item
  Suppose that \(\mathbf{A} = \mathbf{K}[\mathbf{X}, \mathbf{Y}]\) , where \(\mathbf{K}\) is a field. Then the maximal ideal \(m\) generated by \(\mathbf{X}\) and \(\mathbf{Y}\) is a torsion-free A-module, but not flat. Indeed, consider the ring \(\mathbf{B} = \mathbf{A}/(\mathbf{Y})\) , which is isomorphic to \(\mathbf{K}[\mathbf{X}]\) , hence an integral domain. The B-module \(m_{(B)}\) is isomorphic to \(m/Ym = (X, Y)/(XY, Y^2)\) , in which the class of \(Y\) is torsion. Consequently, \(m_{(B)}\) is not a flat B-module, hence \(m\) is not flat.
\item
  Suppose A is commutative. Let B be the algebra \(\mathbf{A}[\mathbf{X}_1, \ldots, \mathbf{X}_n] / (\mathbf{P})\) , where P is a nonzero polynomial. For every prime ideal p of A, denote by \(\kappa(p)\) the field of fractions of the integral domain \(\mathbf{A} / p\) , \(\mathbf{E}(p)\) the algebra \(\kappa(p) [\mathbf{X}_1, \ldots, \mathbf{X}_n]\) and \(\mathbf{P}(p)\) the image of P in \(\mathbf{E}(p)\) under the canonical mapping.
\end{enumerate}

One can show that, for B to be a flat A-module, it suffices that \(P(p) \neq 0\) for every prime ideal p of A. If A is an integral domain, this condition is necessary.

* In geometric language, consider the projection \(\pi: \text{Spec}(B) \to \text{Spec}(A)\) . For every \(p \in \text{Spec}(A)\) the fiber \(\pi^{-1}(p)\) is identified with the subvariety \(V_p\) of the affine space \(\mathbf{A}_{\kappa(p)}^n = \text{Spec}(E(p))\) defined by \(P(p)\) , and the set \(F\) of \(p\) for which this subvariety is the entire space (i.e.~for which \(P(p) = 0\) ) is a closed subset of

Spec(A). The preceding condition means that this closed set is empty, in other words that for every p the subvariety \(V_{p}\) is a hypersurface in \(\mathbf{A}_{\kappa(p)}^{n}\) . *

\begin{enumerate}
\def\labelenumi{\arabic{enumi})}
\setcounter{enumi}{10}
\tightlist
\item
  Let S and X be two complex analytic spaces, and \(f: \mathbf{X} \to \mathbf{S}\) a morphism. We say that \(f\) is flat at a point \(x\) of \(\mathbf{X}\) if \(\mathcal{O}_{\mathbf{X},x}\) , considered as \(\mathcal{O}_{\mathbf{S},f(x)}\) -module by means of the homomorphism \(f^{*}: \mathcal{O}_{\mathbf{S},f(x)} \to \mathcal{O}_{\mathbf{X},x}\) , is flat. The set of points of \(\mathbf{X}\) where \(f\) is flat is an open subset of \(\mathbf{X}\) , and the restriction of \(f\) to this open set is an open map. If \(\mathbf{X}\) and \(\mathbf{S}\) are connected finite-dimensional analytic manifolds, \(f\) is flat (at every point of \(\mathbf{X}\) ) if and only if \(f(\mathbf{X})\) is open in \(\mathbf{S}\) and the fibers \(f^{-1}(s)\) , for \(s \in f(\mathbf{X})\) all have the same dimension. *
\end{enumerate}

